%%%%%%%%%%%%%%%%%%%%%%%%%%%%%%%%%%%%%%%%%%%%%%%%%%%%%%%%%%%%%%%%%%%%%%%%
\section{Introduction}\label{sec:intro}
%%%%%%%%%%%%%%%%%%%%%%%%%%%%%%%%%%%%%%%%%%%%%%%%%%%%%%%%%%%%%%%%%%%%%%%%

\subsection{Meanders and the KPZ formula}

A \emph{meander of order $n$} is a closed curve in the plane, the \emph{road}, that crosses a
fixed straight line, the \emph{river}, exactly $2n$ times and never crosses itself, counted up to
deformations that keep the pattern of crossings. The river cuts the road into \emph{arches}, the
pieces of road between consecutive crossings; the arches above the river pair up the crossings
without crossing each other, and so do the arches below. Let $M_n$ be the number of meanders of
order $n$; the first values are $1,2,8,42$. Two meanders of orders $a$ and $b$ can be \emph{glued
end to end}: place them side by side on one river, remove the upper arch at the last crossing of
the left one and the upper arch at the first crossing of the right one, and reconnect the four
loose ends by a short arch joining these two crossings and a long arch over it. The two roads merge
into one, giving a meander of order $a+b$. Gluing is injective, since the two pieces are recovered
by undoing the reconnection, so $M_{a+b}\ge M_aM_b$, and $\mu=\lim_n M_n^{1/n}=\sup_nM_n^{1/n}$ exists by Fekete's lemma;
numerically $\mu\approx12.26$~\cite{JG}. The leftover factor $u_n=M_n/\mu^n$ is believed
to decay like a power of $n$:
\begin{equation}\label{eq:conj}
M_n\;\approx\;A\,\mu^n n^{-\alpha},\qquad \alpha=\astar=\frac{29+\sqrt{145}}{12}\approx3.4201 .
\end{equation}
This value was predicted by Di~Francesco, Golinelli and Guitter~\cite{DGG}, and numerical
estimates agree with it~\cite{JG}. It has not been proved.

The prediction comes from physics. A large meander is viewed as a random surface carrying two
curves, river and road, that each visit every part of it. The discrete random surfaces here are
\emph{planar maps}: connected graphs drawn on the sphere without edge crossings, counted up to
deformation. Each kind of decoration drawn on a random surface (a spanning tree, a curve, a spin
configuration) has a \emph{central charge} $c$, a number that measures how strongly the decorated
surface prefers round shapes to shapes with thin necks: $c=0$ for a bare surface, $c=-2$ for a
spanning tree and, it is predicted, for a curve through every point. Independent decorations are
expected to add their central charges, so a meander is predicted to have $c=-4$. The
\emph{KPZ formula} of Knizhnik, Polyakov and Zamolodchikov~\cite{KPZ} turns $c$ into an exponent:
\begin{equation}\label{eq:kpz}
\alpha=1+x,\qquad x+\frac1x=\frac{13-c}{6},\qquad x\ge1 .
\end{equation}
Here $x$ is the exponent of a tail: if, after the exponential growth is divided out, surfaces of
size exactly $n$ carry total weight about $n^{-1-x}$, then surfaces of size at least $n$ carry total
weight about $n^{-x}$. At $c=0$ the
formula gives $\alpha=5/2$, the exponent of planar maps found by Tutte~\cite{Tutte}; at $c=-2$ it
gives $\alpha=3$, the exponent of maps with a spanning tree found by Mullin~\cite{Mullin}; at $c=-4$
the equation $x+1/x=17/6$ has the root
\[
\xstar=\frac{17+\sqrt{145}}{12}\approx2.4201,\qquad 1+\xstar=\astar .
\]
Paper~A explains this argument from the beginning~\cite{PaperA}. For meanders it rests on two
unproved steps: that each curve counts as $c=-2$, and that the two curves are independent enough
at large scales for their central charges to add. Paper~C collects what is proved about the numbers
$M_n$ themselves~\cite{PaperC}.

This paper concerns two continuum settings in which parts of the prediction can be proved. In
both, the random surface is not glued from finitely many polygons but is a fixed square or sphere
whose geometry is distorted by a random function, the Gaussian free field. Each object is explained
when it first appears. The letter $n$ always denotes a size: the order of a meander, or the number
of pieces of a surface.

\subsection{What is new}

In the model of Ang, Park, Pfeffer and Sheffield~\cite{APPS}, the \emph{APPS model}, a square is
repeatedly cut into four wherever the random field is high, with a parameter $Q\ge2$ in the
threshold, and the number $N$ of final squares is the size. The continuum analog of counting
spanning trees is the product of the eigenvalues of the Laplacian, suitably renormalized. APPS
proved that, at each fixed size, weighting by a power of this determinant is the same as changing
$Q$. What was missing for an exponent was the law of $N$, and the comparison across sizes.
\begin{itemize}
\item \textbf{The tail law} (Theorem~\ref{thm:tail}): $\Prob(N\ge n)=n^{-4/\gamma^2+o(1)}$ for every
$Q\ge2$, where $Q=2/\gamma+\gamma/2$, including the critical endpoint $Q=2$.
\item \textbf{The normalization} (Propositions~\ref{prop:det}, \ref{prop:identity}): the determinant
carries a factor proportional to $N^{1/3}$, so its $k$-th power needs the compensation $N^{-k/3}$, as
well as an exponential factor that changes the growth rate but not the exponent.
\item \textbf{The KPZ curve} (Theorem~\ref{thm:kpz}): for every real $k\ge-1/2$, the compensated
$k$-th determinant weight of sizes at least $n$ is $n^{-x_k+o(1)}$ with $x_k+1/x_k=(13+2k)/6$. This
is the KPZ formula~\eqref{eq:kpz} at $c=-2k$; at $k=2$, the analog of two spanning trees,
$x_2=\xstar$, which corresponds to $\astar=1+x_2$. This is a statement about sizes at least $n$. For sizes
exactly $n$, where the exponent would be $\astar$ itself, only the upper bound is proved (Section~\ref{sec:tail}).
\item \textbf{Area is not size} (Proposition~\ref{prop:area}): the field also defines a random area
of the square; given a very large area, the probability that $N=1$ stays bounded away from zero.
\item \textbf{A branching random walk} (Theorem~\ref{thm:brwlocal}): in a simplified version where
the field is replaced by independent Gaussian steps on a tree, whose tail is due to A\"id\'ekon, Hu
and Zindy~\cite{AHZ}, $\Prob(\text{size}=n)$ is between constant multiples of $n^{-1-x}$.
\end{itemize}
The second setting is a random surface with two independent \emph{space-filling} curves, curves
that pass through every point; Borga, Gwynne and Sun conjectured it to be the limit of large random
meanders~\cite{BGS}.
\begin{itemize}
\item \textbf{Near-splits} (Theorem~\ref{thm:nearsplit}): a \emph{near-split with tolerance
$\delta$} means that the region the river covers in the first half of its time and the region the
road covers in some half of its time agree up to area $\delta$. Its probability in the limit object
is between $c_1\delta^{\astar}$ and $C_1\delta^{\xstar+5/7}$, where $\xstar+5/7\approx3.1344$; the
exponent of the lower bound is the predicted one. For finite meanders the same event at tolerance
below $1/(4n)$ is exactly a \emph{thin neck} at the middle of the river: a closed curve meeting the
road only twice that cuts the meander into two equal halves.
\item \textbf{Two glued disks} (Proposition~\ref{prop:disks}): two random disk-shaped surfaces glued
along a common boundary of length about $\sqrt\delta$, given total area one, split evenly within
$\delta$ with probability of order $\delta^{1+4/\gamma^2}$, which is $\delta^{\astar}$ at the meander
value of $\gamma$.
\item \textbf{Conditional transfer} (Propositions~\ref{prop:condlower}, \ref{prop:condupper}): if
near-splits of finite meanders, at tolerances just above one lattice cell, are not much rarer than in the
limit object, then $M_n\ge\mu^nn^{-\astar-o(1)}$; if finite meanders can be coupled with the limit
object with polynomially small errors, then $M_n\le C\mu^nn^{-\alpha'}$ for an explicit exponent
$\alpha'$.
\end{itemize}

\subsection{Status of the results}

The results listed above are new, except those attributed to others. The background in
Section~\ref{sec:model}, the first part of Section~\ref{sec:tilt} and Section~\ref{sec:limit} is
from the literature. Each new result has a written proof in the project repository~\cite{Repo}, and
each was checked by at least one separate, independent reading of that proof;
Appendix~\ref{app:repo} lists the files. Nothing has been reviewed outside the project. A full
independent review of the complete arguments, and a systematic search of the literature for
prior or overlapping results, are still pending. The results we use from the literature were
checked against the arXiv versions of the papers cited, not against the journal versions; in one
place (Section~\ref{sec:tilt}) the arXiv text has a misprint that we work around. Short arguments
are given in full; longer ones are sketched, and their complete case analysis stays in the
repository.

The results are of three kinds, and it matters which is which.
\begin{itemize}
\item \emph{Theorems about the APPS model} and its branching-walk simplification
(Sections~\ref{sec:tilt}--\ref{sec:brw}). The model is well defined and contains no meanders; the
determinant only imitates spanning trees.
\item \emph{Theorems about the conjectured limit} (Theorem~\ref{thm:nearsplit}) and about glued
disks (Proposition~\ref{prop:disks}). These are theorems about well-defined random objects. Their
relevance to meanders rests on the conjecture of Borga, Gwynne and Sun, which is open.
\item \emph{Conditional statements} (Propositions~\ref{prop:condlower}, \ref{prop:condupper}). They
say what the missing steps would give, and assume those steps.
\end{itemize}
The only unconditional statements about finite meanders are the combinatorial facts of
Proposition~\ref{prop:micro} and Lemma~\ref{lem:transfer}. In particular, nothing in this paper is
yet a theorem about the exponent $\alpha$ of finite meanders.

Sections~\ref{sec:model}--\ref{sec:brw} treat the APPS model, Sections~\ref{sec:limit}--\ref{sec:disks}
the conjectured meander limit, Section~\ref{sec:bridge} the transfer to finite meanders, and
Section~\ref{sec:open} open problems.
